\vspace{-4mm}
\section{Existing LVLMs Hallucination Evaluation Benchmarks and Metrics}
\vspace{-2mm}

As shown in Table~\ref{tab:evaluation_frameworks}, existing studies \cite{li-etal-2023-evaluating, Wang2023EvaluationAA,Zhai2023HallESwitchRA, Lovenia2023NegativeOP,Villa2023BehindTM,Petryk2024ALOHaAN,Kaul2024THRONEAO} have primarily focused on \textit{object-level} hallucination, with only a few recent studies~\cite{faithscore, wang2023llm,Jiang2024HalEvalAU,Zhang2024QuantityMT} recognizing the importance of extending hallucinations to other dimensions. Our benchmark \data~ covers hallucination evaluations of \textit{objects}, \textit{attributes}, and \textit{relations}, and we further detail attributes to color and counting, and relations to positional and comparative, to provide a comprehensive and fine-grained evaluation benchmark. 

Regarding benchmark annotations, many existing benchmarks employ different ways of annotating the evaluation datasets \textit{automatically}. For example, \citet{li-etal-2023-evaluating} employ object detectors to identify all objects in an image; \citet{Zhai2023HallESwitchRA} employ GPT-4V(ision) to generate ground-truth annotations. There are also approaches to developing models specifically for automatic evaluation, thereby bypassing the need for benchmark collections process~\cite{Wang2023EvaluationAA, Gunjal2023DetectingAP}. Given the challenges and potential inaccuracies associated with automated models, our study opts to annotate the evaluation dataset \textit{manually} to ensure the annotation accuracy and encompass the distinct categories of hallucinations. 

Additionally, most existing benchmarks focus exclusively on hallucination evaluation, which can favor precise but uninformative model outputs, overlooking the aspect of coverage. To address the issue, we incorporate coverage scores in our evaluation. We note that two relevant concurrent works~\cite{wang2023llm,Zhai2023HallESwitchRA} also include the coverage scores. However, compared with our work, they are either limited in scope, focusing only on objects or simple attributes and relations, or are unable to be adopted in open-vocabulary generation settings. Besides, along with the benchmark, we propose an evaluation metric generalizing their adopted CHAIR metric.